\documentclass[11pt,a4paper]{article}

% ══════════════════════════════════════════════════════════════════════════════
%                              PACKAGES
% ══════════════════════════════════════════════════════════════════════════════

\usepackage[utf8]{inputenc}
\usepackage[T1]{fontenc}
\usepackage[french]{babel}
\usepackage{geometry}
\usepackage{graphicx}
\usepackage{xcolor}
\usepackage{titlesec}
\usepackage{fancyhdr}
\usepackage{hyperref}
\usepackage{booktabs}
\usepackage{tabularx}
\usepackage{multicol}
\usepackage{enumitem}
\usepackage{tikz}
\usepackage{fontawesome5}
\usepackage{tcolorbox}
\usepackage{array}
\usepackage{colortbl}

% ══════════════════════════════════════════════════════════════════════════════
%                              CONFIGURATION
% ══════════════════════════════════════════════════════════════════════════════

\geometry{margin=2cm, top=2.5cm, bottom=2.5cm}

% Couleurs
\definecolor{primary}{RGB}{37, 99, 235}
\definecolor{secondary}{RGB}{16, 185, 129}
\definecolor{dark}{RGB}{31, 41, 55}
\definecolor{light}{RGB}{243, 244, 246}
\definecolor{warning}{RGB}{245, 158, 11}
\definecolor{accent}{RGB}{139, 92, 246}

% Hyperlinks
\hypersetup{
    colorlinks=true,
    linkcolor=primary,
    urlcolor=primary,
    citecolor=primary
}

% Titres
\titleformat{\section}{\Large\bfseries\color{primary}}{\thesection}{1em}{}[\titlerule]
\titleformat{\subsection}{\large\bfseries\color{dark}}{\thesubsection}{1em}{}
\titleformat{\subsubsection}{\normalsize\bfseries\color{dark}}{\thesubsubsection}{1em}{}

% En-têtes
\pagestyle{fancy}
\fancyhf{}
\fancyhead[L]{\small\color{gray}GeoRetail MVP - Executive Summary}
\fancyhead[R]{\small\color{gray}Janvier 2025}
\fancyfoot[C]{\thepage}
\renewcommand{\headrulewidth}{0.4pt}

% Boîtes colorées
\tcbuselibrary{skins,breakable}

\newtcolorbox{highlightbox}[1][]{
    colback=primary!5,
    colframe=primary,
    fonttitle=\bfseries,
    title=#1,
    rounded corners,
    boxrule=0.5pt
}

\newtcolorbox{successbox}[1][]{
    colback=secondary!5,
    colframe=secondary,
    fonttitle=\bfseries,
    title=#1,
    rounded corners,
    boxrule=0.5pt
}

\newtcolorbox{warningbox}[1][]{
    colback=warning!5,
    colframe=warning,
    fonttitle=\bfseries,
    title=#1,
    rounded corners,
    boxrule=0.5pt
}

% ══════════════════════════════════════════════════════════════════════════════
%                              DOCUMENT
% ══════════════════════════════════════════════════════════════════════════════

\begin{document}

% ─────────────────────────────────────────────────────────────────────────────
%                              PAGE DE TITRE
% ─────────────────────────────────────────────────────────────────────────────

\begin{titlepage}
    \centering
    \vspace*{2cm}
    
    % Logo/Icône
    {\Huge\color{primary}\faMapMarkedAlt}
    
    \vspace{1cm}
    
    {\Huge\bfseries\color{dark} GeoRetail MVP}
    
    \vspace{0.5cm}
    
    {\Large\color{gray} Plateforme de Cartographie des Points de Vente au Maroc}
    
    \vspace{2cm}
    
    \begin{tcolorbox}[
        colback=light,
        colframe=primary,
        width=0.8\textwidth,
        arc=3mm,
        boxrule=1pt
    ]
        \centering
        {\Large\bfseries Executive Summary}\\[0.5cm]
        {\large Version 2.0 Final}\\[0.3cm]
        {\normalsize Projet de Fin d'Études - Master DSIA}
    \end{tcolorbox}
    
    \vspace{2cm}
    
    % Statistiques clés
    \begin{tabular}{ccc}
        \begin{tcolorbox}[
            colback=primary,
            colframe=primary,
            width=3.5cm,
            arc=2mm,
            boxrule=0pt
        ]
            \centering
            {\Huge\bfseries\color{white} 2976}\\
            {\small\color{white} Points de Vente}
        \end{tcolorbox}
        &
        \begin{tcolorbox}[
            colback=secondary,
            colframe=secondary,
            width=3.5cm,
            arc=2mm,
            boxrule=0pt
        ]
            \centering
            {\Huge\bfseries\color{white} 12}\\
            {\small\color{white} Régions}
        \end{tcolorbox}
        &
        \begin{tcolorbox}[
            colback=accent,
            colframe=accent,
            width=3.5cm,
            arc=2mm,
            boxrule=0pt
        ]
            \centering
            {\Huge\bfseries\color{white} 100\%}\\
            {\small\color{white} Open Source}
        \end{tcolorbox}
    \end{tabular}
    
    \vfill
    
    {\large\color{gray} Janvier 2025}
    
\end{titlepage}

% ─────────────────────────────────────────────────────────────────────────────
%                              CONTENU
% ─────────────────────────────────────────────────────────────────────────────

\section{Résumé Exécutif}

\begin{highlightbox}[\faLightbulb\ Vision du Projet]
GeoRetail est une plateforme de cartographie intelligente des points de vente au Maroc, développée comme \textbf{Minimum Viable Product (MVP)} démontrant les compétences techniques en développement full-stack, intégration d'APIs géospatiales, et conception d'architecture data-driven.
\end{highlightbox}

\subsection{Objectifs Atteints}

\begin{itemize}[leftmargin=*]
    \item[\color{secondary}\faCheckCircle] \textbf{Collecte automatisée} de données géographiques via OpenStreetMap
    \item[\color{secondary}\faCheckCircle] \textbf{Interface cartographique} interactive avec Leaflet.js et clustering
    \item[\color{secondary}\faCheckCircle] \textbf{API RESTful} complète avec FastAPI et documentation Swagger
    \item[\color{secondary}\faCheckCircle] \textbf{Système de filtrage} multi-critères (type, région, enseigne)
    \item[\color{secondary}\faCheckCircle] \textbf{Export multi-format} (CSV, GeoJSON) pour interopérabilité
    \item[\color{secondary}\faCheckCircle] \textbf{Architecture extensible} prête pour sources premium
\end{itemize}

% ─────────────────────────────────────────────────────────────────────────────

\section{Architecture Technique}

\subsection{Stack Technologique}

\begin{center}
\begin{tabular}{>{\columncolor{light}}l l l}
\toprule
\textbf{Composant} & \textbf{Technologie} & \textbf{Version} \\
\midrule
Backend & FastAPI (Python) & 0.109+ \\
Frontend & HTML5/CSS3/JavaScript & ES6+ \\
Cartographie & Leaflet.js + MarkerCluster & 1.9.4 \\
Données & OpenStreetMap (Overpass API) & - \\
Géocodage & Nominatim & - \\
Base de données & JSON (prototype) & - \\
\bottomrule
\end{tabular}
\end{center}

\subsection{Flux de Données}

\begin{center}
\begin{tikzpicture}[
    node distance=2.5cm,
    box/.style={rectangle, draw=primary, fill=primary!10, rounded corners, minimum width=2.5cm, minimum height=1cm, align=center, font=\small},
    arrow/.style={->, >=stealth, thick, color=dark}
]
    \node[box] (osm) {OpenStreetMap\\Overpass API};
    \node[box, right of=osm] (backend) {FastAPI\\Backend};
    \node[box, right of=backend] (db) {JSON\\Database};
    \node[box, right of=db] (frontend) {Leaflet\\Frontend};
    
    \draw[arrow] (osm) -- (backend);
    \draw[arrow] (backend) -- (db);
    \draw[arrow] (db) -- (frontend);
    \draw[arrow, bend left=30] (frontend) to node[above, font=\tiny] {API Calls} (backend);
\end{tikzpicture}
\end{center}

% ─────────────────────────────────────────────────────────────────────────────

\section{Données Collectées}

\subsection{Statistiques Actuelles}

\begin{center}
\begin{tabular}{ll}
\begin{minipage}{0.45\textwidth}
\begin{successbox}[\faDatabase\ Volume de Données]
\begin{itemize}[leftmargin=*,nosep]
    \item \textbf{Total POS:} 2 976
    \item \textbf{Commerce moderne:} 1 039 (35\%)
    \item \textbf{Commerce traditionnel:} 1 937 (65\%)
    \item \textbf{Source:} 100\% OpenStreetMap
\end{itemize}
\end{successbox}
\end{minipage}
&
\begin{minipage}{0.45\textwidth}
\begin{highlightbox}[\faMapMarkerAlt\ Couverture Géographique]
\begin{itemize}[leftmargin=*,nosep]
    \item \textbf{Fès:} 1 464 POS
    \item \textbf{Casablanca:} 332 POS
    \item \textbf{Tanger:} 282 POS
    \item \textbf{Marrakech:} 276 POS
    \item \textbf{Agadir:} 187 POS
\end{itemize}
\end{highlightbox}
\end{minipage}
\end{tabular}
\end{center}

\subsection{Répartition par Catégorie}

\begin{center}
\begin{tabular}{lrr}
\toprule
\textbf{Catégorie} & \textbf{Nombre} & \textbf{Pourcentage} \\
\midrule
Épicerie & 1 076 & 36.2\% \\
Hypermarché & 793 & 26.6\% \\
Boucherie & 379 & 12.7\% \\
Boulangerie & 376 & 12.6\% \\
Supermarché & 235 & 7.9\% \\
Autres & 117 & 3.9\% \\
\bottomrule
\end{tabular}
\end{center}

% ─────────────────────────────────────────────────────────────────────────────

\section{Sources de Données}

\subsection{Sources Gratuites Implémentées}

\begin{successbox}[\faCheckCircle\ OpenStreetMap - Source Principale]
\begin{itemize}[leftmargin=*,nosep]
    \item \textbf{API:} Overpass API (requêtes géographiques)
    \item \textbf{Couverture:} Environ 3 000 POS au Maroc
    \item \textbf{Avantages:} Gratuit, données collaboratives, mise à jour continue
    \item \textbf{Limitations:} Couverture inégale selon les régions
\end{itemize}
\end{successbox}

\vspace{0.5cm}

\begin{successbox}[\faCheckCircle\ Nominatim - Géocodage Complémentaire]
\begin{itemize}[leftmargin=*,nosep]
    \item \textbf{Fonction:} Recherche par nom d'enseigne et ville
    \item \textbf{Rate limit:} 1 requête/seconde (respecté)
    \item \textbf{Utilisation:} Enrichissement de données existantes
\end{itemize}
\end{successbox}

\subsection{Options Premium Documentées (Non Implémentées)}

\begin{warningbox}[\faExclamationTriangle\ Sources Payantes - Pour Référence]
\begin{center}
\begin{tabular}{lll}
\toprule
\textbf{Source} & \textbf{Couverture Estimée} & \textbf{Coût} \\
\midrule
Google Places API & Environ 300 000 POS & 17 USD/1000 req \\
HERE Places API & Environ 150 000 POS & 5 USD/1000 req \\
Foursquare Places & Environ 100 000 POS & 99 USD/mois \\
\bottomrule
\end{tabular}
\end{center}
\textit{Note: Ces sources nécessitent des budgets ou partenariats commerciaux non disponibles pour un projet académique.}
\end{warningbox}

% ─────────────────────────────────────────────────────────────────────────────

\section{Fonctionnalités de l'Application}

\subsection{API REST Endpoints}

\begin{center}
\begin{tabular}{llp{6cm}}
\toprule
\textbf{Méthode} & \textbf{Endpoint} & \textbf{Description} \\
\midrule
GET & \texttt{/pos} & Liste tous les POS avec filtres \\
POST & \texttt{/pos} & Crée un nouveau POS \\
GET & \texttt{/pos/\{id\}} & Détail d'un POS \\
DELETE & \texttt{/pos/\{id\}} & Supprime un POS \\
GET & \texttt{/pos/search/nearby} & Recherche par proximité \\
GET & \texttt{/pos/search/bbox} & Recherche par zone \\
POST & \texttt{/collect/osm} & Lance collecte OSM \\
POST & \texttt{/import/csv} & Import fichier CSV \\
GET & \texttt{/export/geojson} & Export GeoJSON \\
GET & \texttt{/stats} & Statistiques globales \\
\bottomrule
\end{tabular}
\end{center}

\subsection{Interface Utilisateur}

\begin{multicols}{2}
\textbf{Fonctionnalités Cartographiques:}
\begin{itemize}[leftmargin=*,nosep]
    \item Carte interactive Leaflet
    \item Clustering intelligent des marqueurs
    \item Marqueurs différenciés par type
    \item Popups informatifs
    \item Zoom et navigation fluides
\end{itemize}

\columnbreak

\textbf{Fonctionnalités de Gestion:}
\begin{itemize}[leftmargin=*,nosep]
    \item Filtrage multi-critères
    \item Ajout manuel de POS
    \item Import CSV/JSON
    \item Export des données
    \item Tableau de bord statistiques
\end{itemize}
\end{multicols}

% ─────────────────────────────────────────────────────────────────────────────

\section{Compétences Démontrées}

\begin{highlightbox}[\faGraduationCap\ Compétences Techniques Validées]
\begin{multicols}{2}
\textbf{Développement Backend:}
\begin{itemize}[leftmargin=*,nosep]
    \item API RESTful avec FastAPI
    \item Programmation asynchrone (async/await)
    \item Validation de données (Pydantic)
    \item Gestion d'erreurs HTTP
\end{itemize}

\columnbreak

\textbf{Développement Frontend:}
\begin{itemize}[leftmargin=*,nosep]
    \item JavaScript ES6+
    \item Manipulation DOM
    \item Appels API fetch
    \item UI/UX moderne
\end{itemize}
\end{multicols}

\vspace{0.5cm}

\begin{multicols}{2}
\textbf{Data Engineering:}
\begin{itemize}[leftmargin=*,nosep]
    \item Intégration APIs externes
    \item ETL (Extract, Transform, Load)
    \item Déduplication de données
    \item Calculs géospatiaux
\end{itemize}

\columnbreak

\textbf{Architecture:}
\begin{itemize}[leftmargin=*,nosep]
    \item Architecture modulaire
    \item Séparation des concerns
    \item Documentation API (Swagger)
    \item Gestion de configuration
\end{itemize}
\end{multicols}
\end{highlightbox}

% ─────────────────────────────────────────────────────────────────────────────

\section{Évolutions Potentielles}

\subsection{Court Terme (v2.1)}

\begin{itemize}[leftmargin=*]
    \item Migration vers PostgreSQL/PostGIS pour performances
    \item Authentification utilisateur (JWT)
    \item Cache Redis pour requêtes fréquentes
\end{itemize}

\subsection{Moyen Terme (v3.0)}

\begin{itemize}[leftmargin=*]
    \item Intégration Google Places (avec budget)
    \item Application mobile React Native
    \item Dashboard analytics avancé
    \item API publique avec rate limiting
\end{itemize}

\subsection{Long Terme (Vision)}

\begin{itemize}[leftmargin=*]
    \item Plateforme SaaS multi-tenant
    \item Machine Learning pour prédiction de zones
    \item Partenariats avec distributeurs
    \item Expansion régionale (Afrique du Nord)
\end{itemize}

% ─────────────────────────────────────────────────────────────────────────────

\section{Conclusion}

\begin{highlightbox}[\faStar\ Bilan du Projet]
Le MVP GeoRetail démontre avec succès la capacité à concevoir et développer une plateforme de cartographie commerciale complète, en utilisant exclusivement des technologies et données open source.

\vspace{0.5cm}

\textbf{Points Forts:}
\begin{itemize}[leftmargin=*,nosep]
    \item Architecture technique robuste et extensible
    \item Collecte automatisée fonctionnelle (2 976 POS)
    \item Interface utilisateur moderne et intuitive
    \item Documentation complète (API Swagger)
    \item Code maintenable et bien structuré
\end{itemize}

\vspace{0.5cm}

\textbf{Valeur Ajoutée:}

Le projet transforme les contraintes de ressources (absence de budget pour APIs premium) en démonstration de compétences d'ingénieur: trouver des solutions créatives, maximiser les sources gratuites, et documenter clairement les options d'évolution pour un contexte de production.
\end{highlightbox}

\vfill

\begin{center}
\rule{0.5\textwidth}{0.4pt}

\vspace{0.5cm}

{\small\color{gray}
\textbf{GeoRetail MVP v2.0 Final}\\
Projet de Fin d'Études - Master Data Science et Intelligence Artificielle\\
Janvier 2025
}
\end{center}

\end{document}
